\cleardoublepage
\chapter{Conclusions}
In this work, we addressed the problem of patient specificity mainly
for the context of robotised lung ultrasound. Through a model that
embeds the pose and the shape of the human body, we were able to
generate an anatomical skeleton. We then implemented an automatic way
of generating VF from the skeletal model, and used
them to constrain the physical interaction to the intercostal areas. The framework has been validated through quantitative experiments:
at first, we measured the repeatability and generalisation capability
of the body modelling approach; secondly, we evaluated the impact of
the VF in the examination by monitoring the execution time in the 2 and 4 frontal areas of a clinically validated ultrasound
protocol. Results show higher-than-acceptable body modelling
repeatability, while patient subjectivity is taken is appropriately considered. The ultrasound protocol was executed faster with the
VF, indicating a promising research direction.
Potential future work includes the integration of real-time
measurements, such as contact force, to enhance the body model with
viscoelastic properties, e.g., \cite{beber2024passive}, and to
compensate for the limitations of the vision-based approach. We also
plan to conduct further studies involving multiple clinicians and a
larger pool of subjects.